\chapter{Judges and Officials Rules}

\section{Racing Officials}

\subsection{Race Director}
	\begin{enumerate}[leftmargin=1.5em,labelsep=*]
		\item The Race Director is the head organizer and administrator of road race events.
		The Race Director is the highest authority on everything to do with the road race events, except for decisions on rules and results.
		With the Convention Host, the Race Director determines the course, obtains permits, interfaces with the community, and determines the system used to run the event.
		\item The Race Director is responsible for the logistics, equipment for all road racing events.
		\item With the Referee, the Race Director is in charge of keeping events running on schedule, and answers all questions not pertaining to rules and judging.
		\item The Race Director is responsible for ensuring that a race briefing is held with the riders before each race, in which course-specific aspects, 
        in particular safety aspects, are explained.
	\end{enumerate}

\subsection{Referee}
	\begin{enumerate}[leftmargin=1.5em,labelsep=*]
		\item The Referee is the head racing official. 
		The referee is responsible for making sure that the competition rules are observed and for deciding on all related questions that arise during the event.
		\item The referee has full control and authority over all judges and must instruct them on all details and regulations related to the competition.
		\item The referee must ensure that all necessary judges are on their assigned places.
		The referee may replace absent, incapacitated or inadequate judges with others and may appoint additional judges.
		The referee shall ensure that the judges do not interfere in the competition in a biased manner.
		\item The referee can ban persons from the competition venue for the duration of the competition, if they significantly disturb the execution of the competition.
		\item The referee has the exclusive right to disqualify riders in case of violations of the competition rules, with the exception of violations concerning the start.
		Violations of the competition rules can be determined by their own observations or in reports of the responsible judges.
		Violations of the competition rules have to be reported to the referee with the following information: 
        Position and name of the judge, competition, start number of the rider, as well as a clear description of the violation.
	\end{enumerate}

\subsection{Starter}
	\begin{enumerate}[leftmargin=1.5em,labelsep=*]
		\item The Starter starts races and calls riders back in the event of false starts.
		The Responsibilities follow from the text in 3B.5.2, Starting and 3B.5.3, False Starts.
		\item If a verbal (spoken) count is used, the time between (the start of) each of these elements must be the same, and should be approximately 1 second.
		Starters should practice this before the races begin.
		Timing of the count is very important for an accurate start.
%TODO ---  There are international competitions where participants speak other language than English, for example LAUCC where the common language is Spanish
		This count is to be given in English at Unicon or international competitions.
		At other competitions, English is optional.
		\item The Starter checks riders for correct unicycles and safety equipment and will remove from the starting line-up 
        any riders not properly equipped to race, including riders with dangerously loose shoelaces.
		\item The starter has to take a position for the start from which they have an unobstructed view of the riders and 
        the start command and signal can be easily perceived by the riders.
	\end{enumerate}

\subsection{Timekeeper}
	\begin{enumerate}[leftmargin=1.5em,labelsep=*]
%TODO --- Change "the conformity" to "conforming"
		\item The Timekeeper must be responsible for the conformity to the rules and the functioning of the timing system.
		\item The Timekeeper supervises the timing system and if a fully automatic timing and photo finish system is used, makes sure that the camera is correctly aligned.
		\item The Timekeeper (in conjunction with an adequate number of assistants) must determine the official times of the riders.
		The timekeeper must ensure that these results are correctly entered in or transferred to the competition results system.
	\end{enumerate}

\subsection{Course Judge}
	\begin{enumerate}[leftmargin=1.5em,labelsep=*]
		\item The Course Judge observes the adherence to the corresponding competition rules during the race.
		Sufficient number of Course Judges are to be appointed according to the discipline to guarantee an adequate supervision of the adherence to the rules.
		\item The Course Judge can be assigned the task of securing the race course and directing the riders to the correct way.
	\end{enumerate}

\section{Officials Can Compete}

The Referee may not compete in any competition where they may be required to make a decision.
The Race Director may compete, as long as the race course has been announced early enough that the Race Director does not have an advantage from knowledge of the course.

\section{Consequences of Infractions}

The Referee has final say on whether a rider's safety equipment is sufficient.
The Starter will remove from the starting line-up any riders not properly equipped to race, including riders with dangerously loose shoelaces.

A rider who is forced to dismount due to interference by another rider may file a protest immediately at the end of the race.
Riders who intentionally interfere with other riders may receive from the Referee a warning, a loss of placement (given the next lower finishing place), 
disqualification from that race/event, or suspension from all races.
